\section{Introduction}
\label{sec:introduction}

\textit{Pseudomonas aeruginosa} appears on the World Health Organization's priority list of antibiotic-resistant pathogens because of its capacity to accumulate resistance to nearly every class of antipseudomonal agent \parencite{Tacconelli2018_who}. That framing needs one correction, and it cuts against rather than for the urgency of this review's question: the 2024 update to the priority list \emph{downgraded} carbapenem-resistant \textit{P.\ aeruginosa} from Critical to High \parencite{WHO2024_bppl}. We note it explicitly rather than quietly citing the superseded 2018 list. The organism remains a High-priority pathogen for which salvage options are genuinely scarce, which is sufficient motivation; overstating its rank is not necessary and would not survive review.

A second correction is owed to the premise itself. It is no longer accurate to say the antipseudomonal pipeline has stalled: ceftolozane--tazobactam, ceftazidime--avibactam, imipenem--relebactam and cefiderocol all reached clinical use in the last decade and between them recovered activity against many isolates that older definitions call multidrug-resistant. The scarcity this review addresses is therefore not an absence of new agents but the narrower and harder situation those agents leave behind --- the patient in whom they have also failed, or cannot be used. That distinction has a formal counterpart. The Magiorakos categories (MDR/XDR/PDR) count how many antimicrobial \emph{categories} an isolate resists, a measure the new agents have made less decision-relevant; \textcite{Kadri2018_dtr} instead define \emph{difficult-to-treat resistance} (DTR) as non-susceptibility to all \emph{first-line} agents --- the carbapenem, extended-spectrum cephalosporin, $\beta$-lactam/$\beta$-lactamase-inhibitor and fluoroquinolone categories --- which is the point at which a clinician is forced onto reserve therapy. The reserve agents themselves, which \citeauthor{Kadri2018_dtr} name as aminoglycosides, polymyxins and tigecycline, do \emph{not} enter the determination: an isolate susceptible only to colistin is DTR precisely because colistin is what is left. DTR, not MDR, is the category current guidance uses for treatment decisions in \textit{P.\ aeruginosa}. We retain the Magiorakos axis as primary because it is what the published literature reports, and we carry DTR as a separate stratum precisely to show what that reporting practice costs (Sections~\ref{sec:methods-eligibility} and~\ref{sec:results-synthesis}). Once an isolate meets criteria for multidrug resistance (MDR), extensive drug resistance (XDR), or pandrug resistance (PDR) under the Magiorakos framework \parencite{Magiorakos2012_mdrxdr}, treatment options narrow quickly. The antipseudomonal development pipeline has not kept pace, and bacteriophage therapy---administered alone or alongside antibiotics---has moved from isolated compassionate-use reports toward a small but growing set of early-phase trials as a salvage option for infections that no longer respond to conventional regimens.

This evidence base has already been synthesized many times. Since \textcite{ElHaddad2019_eskape} reviewed phage therapy against ESKAPE pathogens in 2019, at least eight further systematic reviews, meta-analyses, and scoping reviews of human phage therapy have appeared, including a large multi-author synthesis in \textit{Lancet Infectious Diseases} \parencite{Uyttebroek2022_lancetid} and, most recently, an individual-patient-data meta-analysis of 130 studies spanning multiple pathogens \parencite{Liu2025_ijaa}. None restricts its pooling to \textit{P. aeruginosa} specifically, and none stratifies pooled outcomes by resistance category or by route of phage administration---the two variables clinicians most need resolved when choosing a regimen for a given patient and infection site.

Our original aim was narrower still: to compare, within the same studies, phage-antibiotic combination therapy against antibiotic monotherapy for MDR/XDR/PDR \textit{P. aeruginosa} and to estimate a pooled effect size for clinical success. A systematic search of PubMed/MEDLINE and ClinicalTrials.gov found no study reporting both a genuine antibiotic-monotherapy arm and a phage-antibiotic combination arm in a \textit{P. aeruginosa}-specific, MDR/XDR/PDR-defined population, so this comparison could not be constructed under any available within-study design. We therefore reframed the review around pooled proportions rather than a paired effect estimate: the proportion of patients achieving clinical success and the proportion experiencing a safety event across the single-arm and comparative studies that do exist, with pre-specified stratification by resistance category (MDR, XDR, or PDR), route of administration, and treatment modality (antibiotic monotherapy vs.\ phage-antibiotic combination).

Specifically, among patients with MDR, XDR, or PDR \textit{P. aeruginosa} infections treated with bacteriophage therapy, we asked: what proportion achieve clinical success, what proportion experience an adverse event, and do these proportions differ by resistance category, administration route, or treatment modality?

We found the current evidence base able to answer only part of the stratified question it was designed to ask. Pooling across all included studies yielded a clinical-success proportion of \CsOverallEst(\CsOverallCI), directionally consistent with the high success rates reported by prior, non-stratified syntheses---but this figure comes from a compassionate-use population selected for refractory infection after antibiotic failure and should not be read as evidence that phage therapy outperforms, or underperforms, any comparator. Stratification succeeded, but only just, for one resistance class: MDR isolates specifically now have a pooled estimate (\CsResMDREstCI clinical success). We present this as a single-cohort characterization rather than a headline achievement: it draws two-thirds of its patients from one cohort and clears the pre-specified evidentiary threshold only under author-reported resistance labels---restricted to independently re-derived classifications, the MDR cell collapses from \ClassMDRStudiesAll studies and \ClassMDRPatientsAll patients to \ClassMDRStudiesVerified studies and \ClassMDRPatientsVerified patients, and cannot be pooled at all. XDR, PDR, every specific administration route, and modality all fell short of the number of studies and patients needed for a stable pooled estimate, and despite a targeted search, no eligible antibiotic-monotherapy arm was identified---so the MDR-versus-XDR comparison that motivated this review remains unanswered even though MDR alone is not. Documenting exactly which parts of this stratified question the current evidence can and cannot answer, rather than forcing an estimate the data cannot support or dismissing the one part that succeeded, is the contribution of this review.
